\chapter{Introduzione}

La programmazione coreografica consente di descrivere sistemi distribuiti a partire da una specifica globale delle interazioni tra processi. Questo paradigma di programmazione rende esplicita la struttura comunicativa del sistema e mostra il suo comportamento complessivo. Il linguaggio Racing Choreographies estende tale modello introducendo il costrutto di race, che permette di esprimere competizione tra processi. Questa tesi presenta la progettazione e l’implementazione di un parser e di un simulatore per tale linguaggio.

\section{Scopo della tesi}

Lo scopo di questa tesi è realizzare un tool per l’esecuzione del linguaggio Racing Choreographies. Il lavoro consiste nell’implementazione di un parser in grado di riconoscere programmi scritti secondo una grammatica e di un simulatore capace di eseguirli secondo una semantica operazionale.

Il contributo dell’elaborato consiste nella definizione di una pipeline che integra:

\begin{itemize}
    \item analisi lessicale e sintattica;
    \item costruzione di un Abstract Syntax Tree;
    \item validazione statica del programma;
    \item una macchina astratta per l’esecuzione dei programmi.
\end{itemize}

La tesi è prevalentemente implementativa: a partire da una sintassi e da una semantica definite formalmente, l’obiettivo è rendere i programmi del linguaggio concretamente eseguibili.

\section{Contesto}

Una coreografia è una specifica globale che descrive in modo unitario le interazioni tra un insieme finito di processi concorrenti. Essa definisce esplicitamente le comunicazioni, le scelte e il flusso di controllo a livello dell’intero sistema. Il programma risultante rappresenta quindi il comportamento complessivo del sistema come un’unica struttura.

Questo paradigma di programmazione è interessante perché permette di ragionare sul sistema a livello globale prima ancora di considerare l’implementazione dei singoli processi partecipanti e inoltre rende esplicite le dipendenze tra processi e riduce il rischio di errori come mismatch nelle comunicazioni o deadlock dovuti a specifiche incoerenti.

Il linguaggio Racing Choreographies estende il modello coreografico classico: tramite le race si rende esplicito un fenomeno tipico dei sistemi distribuiti: la competizione tra eventi concorrenti e la necessità di determinare un vincitore, mantenendo la possibilità di gestire il processo perdente.

\section{Struttura della tesi}

La tesi è organizzata in sei capitoli, che seguono una progressione dal contesto teorico all’implementazione del tool sviluppato.

Il Capitolo 2 descrive il background teorico.

Il Capitolo 3 presenta la definizione formale del linguaggio dal punto di vista sintattico e semantico.

Il Capitolo 4 presenta l’implementazione del parser.

Il Capitolo 5 presenta l’implementazione del simulatore.